%! Orthogonality of the Four Subspaces — Unit 4, Lesson 4.1 — Exit Ticket
\documentclass[10pt]{article}
\usepackage{linalg-article}
\usepackage{linalg-boxes}
\newcommand{\TallMath}[1]{$\displaystyle #1\rule[-1.4em]{0pt}{3.2em}$}
\newcommand{\xx}{\mathbf{x}}
\newcommand{\T}{\mathsf{T}}

\begin{document}

\pageheader{Unit 4, Lesson 4.1}{Exit Ticket}
\namedateperiod

\vspace{0.15in}

No notes. Show your reasoning.

\begin{enumerate}[leftmargin=1.8em, itemsep=15pt]
  \item \textbf{Check the pairing.} A matrix $A$ has rows $(1,2,3)$ and $(2,1,3)$, and
        $\xx=\begin{bmatrix} 1 \\ 1 \\ -1 \end{bmatrix}$ is in its nullspace. Dot $\xx$ with
        \emph{each} row. What do the two results show about $N(A)$ and the row space?
        \vspace{1.5cm}
  \item \textbf{Count \& name.} An $m\times n$ matrix has $m=3$, $n=5$, rank $r=2$. Give
        $\dim N(A)$ and $\dim C(A^{\T})$, and check they fill $\mathbb{R}^5$. Which two subspaces are
        orthogonal complements inside $\mathbb{R}^5$?
        \writelines{2}
  \item \textbf{Explain.} Two perpendicular \emph{lines} lie in $\mathbb{R}^3$. Are they orthogonal
        complements? Answer and justify in one sentence, using dimensions.
        \writelines{2}
\end{enumerate}

\end{document}
